\section{Anexos}\label{anexos}

Los anexos se numeran con letras. Los anexos B a E (minutas de las Sprint
Reviews, detalle de historias y casos de prueba, manual de uso y aviso de
privacidad) se agregan en el corte que los requiere.

\subsection{Anexo A. Guion y minuta de la entrevista con el cliente}\label{anexo-a}

\subsubsection{A.1 Guion de la entrevista inicial}\label{anexo-a1}

Objetivo: entender cómo opera hoy el taller, qué problemas le cuestan
dinero o clientes y qué reglas sigue aunque no estén escritas. Las
preguntas se agrupan por etapa del servicio; en cada una se pide un
ejemplo reciente y se pregunta por las excepciones (quién puede romper la
regla y si queda registro).

\paragraph{Negocio y operación general}
\begin{enumerate}
\tightlist
\item ¿Qué tipos de vehículos atienden y cuáles no reciben?
\item ¿Cuántas bahías de trabajo tienen y cuántos vehículos caben en
espera? ¿Qué pasa cuando todo está ocupado?
\item ¿Quiénes trabajan en el taller y qué hace cada uno durante el día?
\item ¿Qué problemas les cuestan hoy más dinero o más clientes?
\end{enumerate}

\paragraph{Citas}
\begin{enumerate}
\tightlist
\setcounter{enumi}{4}
\item ¿Cómo agenda hoy un cliente su cita y quién la registra?
\item Si un cliente pidiera la cita en línea, ¿debe apartar lugar de
inmediato o alguien la confirma?
\end{enumerate}

\paragraph{Recepción del vehículo}
\begin{enumerate}
\tightlist
\setcounter{enumi}{6}
\item ¿Qué datos del vehículo y del cliente anotan al recibirlo?
\item ¿Han tenido reclamos por golpes o faltantes? ¿Cómo los resolvieron
y qué evidencia tenían?
\item ¿Qué hacen cuando no funciona el indicador de combustible?
\end{enumerate}

\paragraph{Diagnóstico, bahías y mecánicos}
\begin{enumerate}
\tightlist
\setcounter{enumi}{9}
\item ¿Quién decide a qué bahía y a qué mecánico va cada vehículo?
\item ¿Pueden trabajar varios mecánicos en la misma orden? ¿Quién
responde por el trabajo?
\item Si una refacción tarda varios días en llegar, ¿qué pasa con el
vehículo y con la bahía?
\end{enumerate}

\paragraph{Cotización y autorización}
\begin{enumerate}
\tightlist
\setcounter{enumi}{12}
\item ¿Cómo autoriza hoy el cliente una cotización y qué queda como
prueba?
\item ¿Han tenido disputas por cobros que el cliente dice no haber
autorizado? ¿Cómo se resolvieron?
\item ¿El cliente puede aceptar unos conceptos y rechazar otros?
\item ¿Qué hacen si durante la reparación aparece un trabajo que sube el
costo?
\item ¿Cuánto tiempo es válida una cotización?
\end{enumerate}

\paragraph{Cobro, entrega y seguimiento}
\begin{enumerate}
\tightlist
\setcounter{enumi}{17}
\item ¿Se entrega alguna vez un vehículo sin pagar completo? ¿Quién lo
autoriza y para qué clientes?
\item ¿Cuántas llamadas reciben para preguntar por el avance y quién las
atiende?
\item ¿Durante cuánto tiempo debería el cliente poder consultar su
vehículo en línea?
\item ¿Qué garantía dan sobre los trabajos y cómo la consultan hoy?
\end{enumerate}

\paragraph{Infraestructura}
\begin{enumerate}
\tightlist
\setcounter{enumi}{21}
\item ¿Qué equipos tienen (computadoras, tabletas, internet) y qué tan
seguido se va la luz o el internet?
\end{enumerate}

\subsubsection{A.2 Minuta de la entrevista del 19/10/2026}\label{anexo-a2}

{\small
\begin{longtable}[]{@{}
  >{\raggedright\arraybackslash}p{0.25\linewidth}
  >{\raggedright\arraybackslash}p{0.7\linewidth}@{}}
\toprule
\textbf{Dato} & \textbf{Contenido} \\
\midrule
\endhead
\bottomrule
\endlastfoot
Fecha & 19/10/2026 \\
Hora y modalidad & \pendiente{} \\
Asistentes por el cliente & \pendiente{} \\
Asistentes por el equipo & \pendiente{} \\
Elaboró & Product Owner (Ángel Yael Avendaño Guzmán) \\
Objetivo & Levantar los requisitos iniciales del sistema a partir del
anteproyecto y del guion A.1. \\
\end{longtable}
}

\textbf{Problemas expresados por el cliente.} Disputas por cobros que el
cliente dice no haber autorizado («dijo que sí, pero cobra otra cosa»),
porque las autorizaciones son verbales o por mensaje; reclamos por daños
sin checklist ni fotografías del estado de llegada; muchas llamadas para
preguntar por el avance, que interrumpen a la recepcionista; bahías y
carga de trabajo asignadas de palabra, con vehículos detenidos sin
justificación; e historial de reparaciones y garantías archivado en
carpetas (sección 3.1).

\textbf{Hallazgos y acuerdos.} Cada hallazgo tiene un identificador que
se cita como fuente en las reglas de negocio de la sección 4.7.

{\small
\begin{longtable}[]{@{}
  >{\raggedright\arraybackslash}p{0.07\linewidth}
  >{\raggedright\arraybackslash}p{0.62\linewidth}
  >{\raggedright\arraybackslash}p{0.25\linewidth}@{}}
\toprule
\textbf{ID} & \textbf{Hallazgo o acuerdo} & \textbf{Se refleja en} \\
\midrule
\endhead
\bottomrule
\endlastfoot
E-01 & Toda autorización de cotización, remota o presencial, debe tener
una doble verificación con un código de 6 dígitos. & RN-02 \\
E-02 & El cliente puede aprobar o rechazar por separado los conceptos
opcionales de una cotización; solo se ejecuta lo autorizado. & RN-03 \\
E-03 & El enlace de seguimiento del cliente funciona mientras el vehículo
está en el taller y vence al entregarlo y saldarlo. & RN-01 \\
E-04 & El taller tiene 4 bahías (un vehículo por bahía) y un
estacionamiento de espera para 5 vehículos. & RN-05 \\
E-05 & Si una refacción tarda varios días, el vehículo pasa al
estacionamiento para liberar la bahía. & RN-06 \\
E-06 & Solo se atienden vehículos ligeros: automóviles, motocicletas y
camionetas de hasta 3 toneladas. & RN-07 \\
E-07 & Las citas pedidas en línea no apartan lugar hasta que la
recepcionista las confirma. & RN-08 \\
E-08 & Varios mecánicos pueden colaborar en una orden, pero siempre hay
un responsable. & RN-09 \\
E-09 & No se entrega un vehículo con saldo pendiente; solo el gerente
puede autorizar la entrega a crédito a flotillas empresariales
registradas. & RN-10 \\
E-10 & Hay cortes de luz e internet; el personal necesita seguir
capturando sin conexión. & RNF-08, RNF-09 \\
E-11 & El jefe de mecánicos trabajará con 1 o 2 tabletas en un ambiente
con ruido, polvo y manos sucias. & 4.3.1, RNF-02, RNF-03 \\
\end{longtable}
}

\textbf{Pendientes para la siguiente sesión con el cliente.} Confirmar el
mínimo de 4 fotografías al recibir el vehículo (RN-12) y la fuente de la
garantía de 30 días sobre la mano de obra (RN-13). Responsable: Product
Owner.

\subsection{Anexo F. Atención de observaciones de la Review 1}\label{anexo-f}

Observaciones recibidas el 10/11/2026 sobre la versión v0.1. Cada
observación tiene su issue con las etiquetas \texttt{tipo:defecto-doc},
severidad y \texttt{origen:cliente}, y se corrige en su propia rama con
un Pull Request que cierra el issue.

Cada corrección aumenta en uno el tercer dígito de la versión del
documento: la v0.2 es la entrega de la Review 1 sin correcciones, la
v0.2.1 incluye la primera corrección, la v0.2.2 la segunda, y así
sucesivamente. La versión se registra en la portada, en el control de
versiones, en esta tabla y con una etiqueta de Git del mismo nombre.

{\small
\begin{longtable}[]{@{}
  >{\raggedright\arraybackslash}p{0.06\linewidth}
  >{\raggedright\arraybackslash}p{0.09\linewidth}
  >{\raggedright\arraybackslash}p{0.07\linewidth}
  >{\raggedright\arraybackslash}p{0.06\linewidth}
  >{\raggedright\arraybackslash}p{0.61\linewidth}@{}}
\toprule
\textbf{ID} & \textbf{Versión} & \textbf{Issue} & \textbf{PR} & \textbf{Cómo se resolvió} \\
\midrule
\endhead
\bottomrule
\endlastfoot
B11 & v0.2.1 & \#9 & \#15 & Se retiraron de 3.4 y de la Figura 3.1 al equipo de
desarrollo y al docente evaluador, que no son stakeholders del producto.
El documento se versiona según la tabla 0.2 del machote: la v0.1
(Review 0) contiene solo portada, 1, 2.1--2.4 y 3, y la v0.2 (Review 1)
agrega 4.1--4.7; cada entrega incluye únicamente las secciones de su
corte. \\
B5 & v0.2.2 & \#10 & \#16 & Se cambió la estrategia de ramas a GitHub Flow: una
rama por historia o tarea con la convención
\texttt{tipo/número-descripción}. Se actualizaron CONTRIBUTING.md, el
README y la sección 2.1; las ramas por integrante
(\texttt{rama/<NOMBRE>}) ya no reciben trabajo nuevo. Las correcciones
de esta revisión ya usan la convención
(\texttt{fix/9-version-stakeholders}, \texttt{docs/10-ramas-por-historia},
\texttt{docs/20-anexo-entrevista}). \\
B3 & v0.2.3 & \#20 & \#21 & Se agregó el Anexo A con el guion de la entrevista
(22 preguntas) y la minuta del 19/10/2026, con los hallazgos E-01 a E-11.
Cada regla de 4.7 que proviene de la entrevista cita en su fuente el
hallazgo de la minuta. Los anexos ya no contienen texto del machote. La
sección 4.11 no forma parte de este corte y se actualiza en la Review 2. \\
B4 & v0.2.4 & \#28 & \#61 & Las 30 historias se registraron como issues
(\#31 a \#60) y 4.5 muestra el número de cada una. En 2.3 se
registraron los puntos del Sprint 1 y la ruta a la minuta de la Review 1
(\ruta{docs/minutas/2026-11-10-review-1.md}). En 12.2 se anotaron los
commits, los Pull Requests y las revisiones de cada integrante, con datos
de GitHub. Siguen pendientes las fechas de los sprints, la minuta de la
Review 0 y las retrospectivas. \\
B2 & v0.2.5 & \#29 & \#62 & Primera parte: se elaboró la propuesta de
alcance (\ruta{docs/04-requisitos/propuesta-release-v0.1.0.md}), que
deja en el Release v0.1.0 las 11 historias del flujo núcleo (40 puntos)
y pasa 14 a la Fase 2. Falta negociarla con el cliente y actualizar 3.3,
3.6, 4.5 y el story map; el issue sigue abierto. \\
B6 & v0.2.6 & \#30 & \#63 & Primera parte: se agregó la sección 2.6 con
la declaración de uso de inteligencia artificial. Falta que cada autor
declare el contenido original de las secciones 1 a 4 y realizar la
sesión interna con minuta; el issue sigue abierto. \\
\end{longtable}
}
